\documentclass[11pt]{article}
\usepackage{color}

%----------Packages----------
\usepackage{amsmath}
\usepackage{amssymb}
\usepackage{amsthm}
\usepackage{dsfont}
\usepackage{mathrsfs}
\usepackage{stmaryrd}
\usepackage[all]{xy}
\usepackage[mathcal]{eucal}
\usepackage{verbatim}  %%includes comment environment
\usepackage{fullpage}  %%smaller margins
\usepackage{multicol}
\usepackage{hyperref}
\usepackage{graphicx}

%Packages added to get the code block
\usepackage{listings}
\usepackage{xcolor}

\definecolor{codegreen}{rgb}{0,0.6,0}
\definecolor{codegray}{rgb}{0.5,0.5,0.5}
\definecolor{codepurple}{rgb}{0.58,0,0.82}
\definecolor{backcolour}{rgb}{0.95,0.95,0.92}

\lstdefinestyle{mystyle}{
    backgroundcolor=\color{backcolour},   
    commentstyle=\color{codegreen},
    keywordstyle=\color{magenta},
    numberstyle=\tiny\color{codegray},
    stringstyle=\color{codepurple},
    basicstyle=\ttfamily\footnotesize,
    breakatwhitespace=false,         
    breaklines=true,                 
    captionpos=b,                    
    keepspaces=true,                 
    numbers=left,                    
    numbersep=5pt,                  
    showspaces=false,                
    showstringspaces=false,
    showtabs=false,                  
    tabsize=2
}

\lstset{style=mystyle}


%---------Commands----------
\def\ZZ{\mathbb{Z}}
\def\QQ{\mathbb{Q}}
\def\PP{\mathbb{P}}
\def\NN{\mathbb{N}}
\def\OO{\mathcal{O}}
\def\AA{\mathbb{A}}
\def\BB{\mathbb{B}}
\def\FF{\mathbb{F}}
\def\LL{\mathbb{L}}
\def\RR{\mathbb{R}}
\def\CC{\mathbb{C}}
\def\GG{\mathbb{G}}
\def\lcm{\mathrm{lcm}}
\def\ord{\mathrm{ord}}
\def\rep{\mathrm{Rep}}
\def\dom{\mathrm{Dom}}
\def\ve{\varepsilon}

\newcommand{\dlegendre}[2]{\ensuremath{\left( \frac{#1}{#2} \right) }}
\newcommand{\tlegendre}[2]{\ensuremath{\big( \frac{#1}{#2} \big) }}

%----------Theorems----------
\theoremstyle{definition}
\newtheorem{theorem}{Theorem}[section]
\newtheorem{proposition}[theorem]{Proposition}
\newtheorem{lemma}[theorem]{Lemma}
\newtheorem{corollary}[theorem]{Corollary}
\newtheorem{definition}[theorem]{Definition}
\newtheorem{example}[theorem]{Example}
\newtheorem*{definition*}{Definition}
\newtheorem{nondefinition}[theorem]{Non-Definition}
\newtheorem{problem}{Problem}[section]
\newtheorem*{tip}{Pro Tip}
\newtheorem*{remark}{Remark}
\numberwithin{equation}{subsection}

\newenvironment{solution}
{%
  \vspace{0.5em}%
  \hrule%
  \vspace{0.2em}%
  \noindent\textbf{Solution:}\!\!\!\!\!
  \pushQED{\qed}
}
{%
  \popQED
  \vspace{.2em}
  \hrule%
  \vspace{0.75em}%
}

\newcommand{\separate}{
    \vspace{.2in}
    \hrule
    \vspace{.2in}
}


\newcommand{\head}[1]{
	\begin{center}
		{\large #1}
		\vspace{.05 in}
	\end{center}
	\bigskip 
}

\newcommand{\blue}[1]{\textcolor{blue}{#1}}
\newcommand{\red}[1]{\textcolor{red}{#1}}

%Problem Set Data
%number = 8
%course = 261
%semester = Fall 2025
%path = /Users/thyde/Documents/Teaching/Vassar/Semesters/2025 Fall/Math 261 Fall 2025/Problem Sets/Problem Set 8

\begin{document}
\pagestyle{plain}

%%---  sheet number for theorem counter
\setcounter{section}{0}   

\head{MATH 999, Summer 2026\\ Note 0: Review of Rings and Fields\\
Trevor Hyde}

%HEADER

\begin{center}
        \emph{Mathematical discoveries, like springtime violets in the woods,\\ have their season which no human can hasten or retard.}
    \vspace{.1in}

    Carl Friedrich Gauss
\end{center}
\vspace{.1in}

In these notes we review the fundamentals of rings and fields to prepare ourselves for the dive into Galois theory.
I encourage you to work through the problems, even if you've seen this before.
Eventually, with enough work, this will all become second nature to you, like reciting your name or tying your shoes.
Until then: practice and review.


\begin{definition}
    A \textbf{(commutative) ring} is a set $R$ together with binary operations called \textbf{addition} $+$ and \textbf{multiplication} $\cdot$ which satisfy the following axioms:
    \begin{enumerate}
        \item (Associativity) For all $a, b, c \in R$ we have
        \begin{align*}
            a + (b + c) &= (a + b) + c\\
            a \cdot (b \cdot c) &= (a \cdot b) \cdot c.
        \end{align*}

        \item (Commutativity) For all $a, b \in R$ we have
        \begin{align*}
            a + b &= b + a\\
            a \cdot b &= b \cdot a.
        \end{align*}

        \item (Distributivity) For all $a, b, c\in R$ we have
        \[
            a\cdot (b + c) = (a\cdot b) + (a \cdot c).
        \]

        \item (Identities) There exists elements $0, 1 \in R$ such that for $a \in R$ we have
        \begin{align*}
            a + 0 &= a\\
            a \cdot 1 &= a.
        \end{align*}

        \item (Additive Inverses) For all $a \in R$ there is an element $b \in R$ such that
        \[
            a + b = 0.
        \]
        We define $-a := b$ and call $-a$ the \textbf{additive inverse} or \textbf{negative} of $a$.
    \end{enumerate}
\end{definition}

\begin{definition}
    Let $R$ be a commutative ring.
    We say $a \in R$ is a \textbf{unit} if there exists some $b \in R$ such that $a\cdot b = 1$.
    Let $R^\times$ denote the set of all units in $R$.
\end{definition}

\begin{definition}
    A \textbf{field} is a commutative ring $R$ such that every nonzero element of $R$ is a unit.
\end{definition}

Fields are the most convenient kinds of rings since we can do all four basic operations (addition, subtraction, multiplication, and division) with them.

\begin{example}
    Basic examples of fields include the rational numbers $\QQ$, the real numbers $\RR$, and the complex numbers $\CC$.
\end{example}

Given a field $K$ we can make a new field $K(x)$ comprised of all rational functions with coefficients in $K$.
That is, elements of $K(x)$ are fractions $\frac{f(x)}{g(x)}$ where $f(x), g(x) \in K[x]$ are polynomials and $g(x)$ is not identically zero.
In this way we can make lots of examples of fields which do not consist of numbers.
Note that we can also iterate this construction, adding more variables to get fields like $K(x_1, x_2,\ldots, x_n)$ of all rational functions in $n$ variables with coefficients in $K$.

Next we review a more sophisticated method for constructing fields via quotient rings.

\begin{definition}
    Let $R$ be a commutative ring.
    An \textbf{ideal} $J \subseteq R$ is a subset which is closed under addition and which satisfies $a J \subseteq J$ for all $a \in R$.
\end{definition}

The main point of ideals is that we can \emph{quotient} a ring by them, which means we can ``divide'' our ring by this ideal to create a new ring.

\begin{definition}
    Let $R$ be a commutative ring and let $J \subseteq R$ be an ideal.
    If $a \in R$, let $a + J := \{a + j : j \in J\}$.
    We call $a + J$ an \textbf{(additive) coset} of $J$.
    Let $R/J := \{a + J : a \in R\}$ be the set of all cosets of $J$.
\end{definition}

\begin{problem}
    Let $R$ be a ring and let $J \subseteq R$ be an ideal.
    \begin{enumerate}
        \item \blue{Prove that $a + J = b + J$ if and only if $a - b \in J$.}
        This is the \textbf{coset equality criteria}.
        Learn it, live it, love it.

        \item \blue{Prove that $R/J$ forms a ring with operations defined by
        \begin{align*}
            (a + J) + (b + J) &:= (a + b) + J\\
            (a + J) \cdot (b + J) &:= (a\cdot b) + J.
        \end{align*}
        and identities $0 + J$ and $1 + J$.
        }
    \end{enumerate}
    
\end{problem}

Recall that a function $\phi : R \to S$ between rings is called a \textbf{homomorphism} if for all $a, b \in R$ we have $\phi(a + b) = \phi(a) + \phi(b)$ and $\phi(a\cdot b) = \phi(a)\cdot \phi(b)$, and $\phi(1) = 1$.

\begin{problem}
    Let $R$ be a ring and let $J \subseteq R$ be an ideal.
    \blue{Prove that $\rho_J : R \to R/J$ defined by $\rho_J(a) := a + J$ is a surjective homomorphism.}
    We call $\rho_J$ the \textbf{reduction modulo $J$} homomorphism or \textbf{quotient by $J$} homomorphism.
\end{problem}

Here is some useful notation borrowed from number theory.
If $R$ is a ring, $a, b \in R$ and $J \subseteq R$ is an ideal, then we say $a$ is \textbf{congruent} to $b$ \textbf{modulo} $J$ and write $a \equiv b \bmod J$ if $a - b \in J$, which is equivalent to $a + J = b + J$.
We can do arithmetic with congruences just like we do with modular arithmetic.
You can think of the ring $R/J$ as the ring of ``modular arithmetic modulo $J$''.
The coset $a + J$ consists precisely of all the elements in $R$ which are congruent to $a$ modulo $J$.

To get some grounding with this notation, note that if $R = \ZZ$ is the ring of integers and $J = \langle m \rangle = \{ma : a \in \ZZ\}$ is the ideal of all multiples of $m$, then $\ZZ/\langle m \rangle$ is the ring of integers modulo $m$ and the coset $a + \langle m\rangle = \{a + mb : b \in \ZZ\}$ is the set of all integers which are congruent to $a$ modulo $m$.
\footnote{An ideal $J \subseteq R$ is said to be \textbf{principal} when $J = \langle m\rangle := \{ma : a \in R\}$ for some $m \in R$.
That is, $J$ consists of all multiples of the element $m$.
}
Note that $a - b \in \langle m \rangle$ if and only if $m$ divides $a - b$.
Hence this fancy definition of $a \equiv b \bmod \langle m \rangle$ is equivalent to the definition of $a \equiv b \bmod m$ we learned in number theory.

You can think of $a - b \in J$ as meaning ``$J$ divides the difference $a - b$'' even though when $J$ is not a principal ideal $\langle m \rangle$ this is not literally true.
Nevertheless, it still remains spiritually true and its a good intuition.

Quotienting a ring by an ideal is a flexible technique for constructing interesting rings.
In particular, we can use this method to construct fields.
A natural question to ask is when will the quotient ring $R/J$ be a field?
This turns out to depend entirely on a special property of the ideal $J$.

\begin{definition}
    We say an ideal $J$ is \textbf{maximal} if $1 \notin J$ and if for all ideals $I$ which satisfy $J \subseteq I \subseteq R$, we either have $I = J$ or $I = R$.    
    That is, an ideal $J \subseteq R$ is maximal if it is as large as it could possibly be without being all of $R$.
\end{definition}

Here is the fundamental result:

\begin{theorem}
\label{thm: maximal field}
    Let $R$ be a commutative ring and let $J \subseteq R$ be an ideal.
    Then $R/J$ is a field if and only if $J$ is a maximal ideal.
\end{theorem}

\begin{proof}
    I'll guide you through the proof.
    First suppose that $R/J$ is a field.
    We want to show that $J$ is maximal.
    So we suppose we have some ideal $I$ such that $J \subseteq I \subseteq R$ and we want to show that either $I = J$ or $I = R$.
    If $I = J$, then we're done.
    Suppose that $I \neq J$.
    Then there exists an element $a \in I \setminus J$.

    \begin{problem}
    \,
        \begin{enumerate}
            \item If $J_1, J_2 \subseteq R$ are ideals, let $J_1 + J_2 := \{j_1 + j_2 : j_i \in J_i\}$.
            \blue{Prove that $J_1 + J_2$ is an ideal containing both $J_1$ and $J_2$.}

            \item Going back to our current setup in the proof, \blue{show that $J \subset J + \langle a \rangle \subseteq I$.}

            \item Note that $a \notin J$ implies that $a \not\equiv 0 \bmod J$.
            \blue{What does $R/J$ being a field allow us to conclude?}

            \item \blue{Unpack definitions to show that $1 \in J + \langle a\rangle \subseteq I$}

            \item \blue{Conclude that $I = R$.}
        \end{enumerate}
    \end{problem}
    That finishes the first half of the proof.

    Now for the reverse implication, suppose that $J$ is a maximal ideal.
    We want to argue that $R/J$ is a field.
    So suppose that $a \in R$ is an element such that $a \not\equiv 0 \bmod J$.
    Then we need to show that $a \bmod J$ is a unit.

    \begin{problem}
    \,
        \begin{enumerate}
            \item \blue{Argue that $J + \langle a \rangle = R$.}
            \item \blue{Unpack what it means to say that $1 \in J + \langle a \rangle$.}
            \item \blue{Conclude that $a \bmod J$ is a unit.}
        \end{enumerate}
    \end{problem}
    We win!
\end{proof}

Let's see this theorem in action.
First consider the ring of integers $\ZZ$.
Division with remainder implies that every ideal in $\ZZ$ is principal (see the Appendix). 
So when is $\langle m \rangle$ a maximal ideal?
The containment $\langle m\rangle \subseteq \langle n \rangle$ is equivalent to $n \mid m$.
So $\langle m \rangle$ will be maximal precisely when $m$ has no proper factors other than $\pm 1$ and $\pm m$, which is equivalent to $m$ being prime.
Hence $\langle m \rangle$ is a maximal ideal if and only if $m$ is a prime number.

If $p$ is a prime, then Theorem \ref{thm: maximal field} implies that $\FF_p := \ZZ/\langle p\rangle$ is a field.
Thus for every prime $p$ there is a field $\FF_p$ with $p$ elements, namely integers modulo $p$.
We proved this fact in number theory without using these terms or these sophisticated tools.
Here we get it as a consequence of more robust machinery.

We can generalize this example to construct fields that are essential for Galois theory.
Let $K$ be a field and let $K[x]$ be the ring of polynomials with coefficients in $K$.
Division with remainder for polynomials implies that every ideal in $K[x]$ is principal (see Appendix).
So when is $\langle m(x)\rangle$ a maximal ideal?
The containment $\langle m(x)\rangle \subseteq \langle n(x)\rangle$ implies that $n(x) \mid m(x)$, meaning that $n(x)$ is a factor of $m(x)$.
Hence $\langle m(x)\rangle$ maximal is equivalent to $m(x)$ having no non-constant factors in $K[x]$, which is to say that $m(x)$ is an \textbf{irreducible} polynomial.
\footnote{A polynomial $m(x) \in K[x]$ being irreducible is not an intrinsic quality of $m$ so much as it is a property of the relationship between $m(x)$ and the coefficient field $K$.
For example, the polynomial $x^2 + 1$ is irreducible as an element of $\RR[x]$, but viewed as an element of $\CC[x]$ it is not irreducible since $x^2 + 1 = (x - i)(x + i)$.}
Hence we have proved the following proposition.

\begin{proposition}
    Let $K$ be a field and let $m(x) \in K[x]$ be an irreducible polynomial, then $K[x]/\langle m(x)\rangle$ is a field.
\end{proposition}

So how should we think about the fields $K[x]/\langle m(x)\rangle$?
I think about this as a two step process:
\begin{enumerate}
    \item First we took the field $K$ and we formally added a new element $x$ to it to create the ring $K[x]$.
    Here $x$ is a blank slate--it satisfies no algebraic identities in $K[x]$.
    Hence $x$ could potentially be \emph{anything}, which is why we often think of it as a variable.

    \item Next we quotient by the ideal $\langle m(x)\rangle$.
    This process has the affect of imposing the algebraic equation ``$m(x) = 0$'', which to use precise notation looks like the congruence $m(x) \equiv 0 \bmod \langle m(x)\rangle$. 
    We then get all algebraic consequences of this identity, which is what the rest of the ideal $\langle m(x)\rangle$ captures.
    After we form this quotient, $x$ is no longer just a variable: $x$ satisfies the polynomial identity $m(x) \equiv 0 \bmod \langle m(x) \rangle$ which means that $x$ is a root of the polynomial $m$ in $K[x]/\langle m(x)\rangle$.
    In other words, we have formally created a field $L := K[x]/\langle m(x)\rangle$ containing $K$ which contains a root of the irreducible polynomial $m(x)$
\end{enumerate}

Even if this construction seems weird and abstract, I claim you've already seen disguised versions of this.
For example, when introducing the complex numbers $\CC$ we typically describe this field as $\CC = \{a + bi : a, b \in \RR,\,  i^2 = -1\}.$
Here $i = \sqrt{-1}$ is introduced as a square root of $-1$, which cannot be a real number and is simply asserted to exist.
Historically people were weirded out by this construction.
Why were we allowed to simply invent numbers that do not exist and imbue them with desirable algebraic properties?
The formal resolution of this uneasiness is the theory of quotient rings that we just reviewed.
To construct the complex numbers we start with the ring $\RR[x]$ and the irreducible polynomial $x^2 + 1 \in \RR[x]$.
Then Theorem \ref{thm: maximal field} implies that $\RR[x]/\langle x^2 + 1\rangle$ is a field and the image of the variable $x$ under the quotient map is precisely a square root of $-1$.
Hence we have an isomorphism $\RR[x]/\langle x^2 + 1\rangle \cong \CC$.
Since $x^2 + 1$ has degree 2, polynomial division with remainder implies that every $f(x) \in \RR[x]$ is congruent to a unique linear polynomial $a + bx$ modulo $x^2 + 1$.
Thus $f(x) = q(x)(x^2 + 1) + (a + bx)$.
``Evaluating'' at $x = i$ this reduces to $f(i) = a + bi$, giving us the standard normal form for complex numbers.

To bring this full-circle, we need to recall the first isomorphism theorem for rings.
Recall that the \textbf{kernel} of a homomorphism $\phi : R \to S$ is $\ker\phi := \{a \in R : \phi(a) = 0\}$.
Kernels are always ideals.

\begin{theorem}[First Isomorphism Theorem]
    If $\phi : R \to S$ is a surjective homomorphism between rings, then the map $\bar\phi : R/\ker\phi \to S$ defined by $\bar\phi(a + \ker\phi) = \phi(a)$ is a well-defined ring isomorphism.
\end{theorem}

\begin{problem}
    \blue{Review the details.}
    The meat of this theorem is showing that the map $\bar\phi$ is well-defined (which is not so much hard as it is tedious).
    Once we know that, everything else is almost immediate.
\end{problem}

If $K \subseteq L$ are fields we say that $L$ is an \textbf{extension} of $K$, and, despite the obvious risk for confusion, we write $L/K$ as a shorthand to mean that $L$ is an extension of $K$.
Note that a subfield of a field is never an ideal, so $L/K$ does not actually parse as a quotient ring.

The ring of integers $\ZZ$ have the special property that for any ring $R$ there is a unique homomorphism $\iota_R : \ZZ \to \RR$.
This is because the definition of a homomorphism forces $\iota_R(1) = 1$ and every integer is constructed by adding 1 to itself some number of times (or taking the negative of such a number) hence all the values of $\iota_R$ are predetermined without us making any choices.
Because $\ZZ$ has this property we say that $\ZZ$ is the \textbf{initial} commutative ring.

\begin{problem}
    Let $K$ be a field.
    \begin{enumerate}
        \item If $\ker\iota_K = \{0\}$, then we say that $K$ has \textbf{characteristic} 0.
        \blue{Prove that in this case $K$ contains a unique subfield isomorphic to $\QQ$.}

        \item Now suppose that $\ker\iota_K = \langle m \rangle$ for some nonzero $m$.
        \blue{Prove that $m$ must be prime.}

        \item \blue{In this case, suppose that $m = p$ is prime and prove that $K$ contains a unique subfield isomorphic to $\FF_p$.}
        We say that $K$ has \textbf{characteristic $p$}.
    \end{enumerate}
    The characteristic is an important invariant which helps us classify fields.
    This problem shows that \emph{every} field is (essentially) an extension of either $\QQ$ or $\FF_p$ for some prime $p$
\end{problem}

It's really easy to construct homomorphisms from polynomial rings.
Let $K$ be a field and suppose that $L \supseteq K$ is some field containing $K$.
If $\alpha \in L$, let $\ve_\alpha : K[x] \to L$ be the \textbf{evaluation at $\alpha$} map defined by $\ve_\alpha(f(x)) := f(\alpha)$.
The map $\ve_\alpha$ is a ring homomorphism (check yourself, the definitions perfectly align to make this immediate).
We define $K[\alpha] \subseteq L$ to be the image of this map.
The kernel of $\ve_\alpha$ is an ideal in $K[x]$.
If $\ker \ve_\alpha = \{0\}$, then that means that $\alpha$ is not the root of \emph{any} polynomial with coefficients in $K$.
In this case we say $\alpha$ is \textbf{transcendental} over $K$.
Otherwise we say $\alpha$ is \textbf{algebraic} over $K$ and $\ker \ve_\alpha = \langle m_\alpha(x)\rangle$ for some unique monic nonzero polynomial $m_\alpha(x) \in K[x]$.
We call $m_\alpha(x)$ the \textbf{minimal polynomial} of $\alpha$ over $K$.
If $\alpha$ is transcendental over $K$, then we define $m_\alpha(x)$ to be the zero polynomial.
Thus the first isomorphism theorem implies that
\[
    K[\alpha] \cong K[x]/\langle m_\alpha(x)\rangle
\]
for all $\alpha \in L$ where $L$ is a field containing $K$.

For example, consider the field $\CC(t)$ where $t$ is a variable.
This field contains $\RR$.
The element $t \in \CC(t)$ is transcendental over $\RR$ since $f(t) \neq 0$ for any nonzero polynomial $f(x) \in \RR[x]$.
The element $i \in \CC(t)$ is algebraic over $\RR$ and the minimal polynomial of $i$ is $m_i(x) = x^2 + 1$.
Hence the first isomorphism theorem implies that
\[
    \CC = \RR[i] \cong \RR[x]/\langle x^2 + 1\rangle.
\]

\begin{problem}
    Let $L/K$ be a field extension.
    Suppose that $\alpha \in L$ is algebraic over $K$.
    \blue{Prove that $m_\alpha(x)$ is irreducible over $K$.}
    Hence it follows that $K[\alpha] \cong K[x]/\langle m_\alpha(x)\rangle$ is a field.
    In this case we can write $K[\alpha]$ or $K(\alpha)$ and they both refer to the same thing: the smallest field extension of $K$ containing $\alpha$.
\end{problem}

\section*{Appendix}

Let's review division with remainder and how it relates to principal ideals.
First, the integers.

\begin{theorem}
    Let $a, b \in \ZZ$ with $b \neq 0$.
    Then there exists unique integers $q, r$ such that
    \begin{enumerate}
        \item $a = qb + r$\\
        \item $0 \leq r < |b|$.
    \end{enumerate}
\end{theorem}

\begin{proof}
    Let $R = \{a - qb : q \in \ZZ \text{ and } a - qb \geq 0\}$.
    This is a nonempty set of natural numbers, hence has a minimal element $r$ by the well-ordering principle.
    Hence there is some integer $q$ such that $a - qb = r$, or equivalently $a = qb + r$.
    
    Without loss of generality we may assume that $b > 0$.
    Suppose that $r \geq b$.
    Then $a - (q + 1)b = a -qb - b = r - b \geq 0$ and $b > 0$ implies that $r - b < r$.
    This contradicts the minimality of $r$.
    Hence $0 \leq r < b$.
    That proves the existence.

    For uniqueness, suppose that $(q_1, r_1)$ and $(q_2, r_2)$ are two pairs satisfying these properties.
    Suppose without loss of generality that $r_1 \geq r_2$.
    Then
    \[
        q_1 b + r_1 = a = q_2 b + r_2
    \]
    which implies that
    \[
        r_1 - r_2 = (q_2 - q_1)b.
    \]
    Thus $b$ divides $r_1 - r_2$ which satisfies $0 \leq r_1 - r_2 < b$.
    The only multiple of $b$ in this range is 0, hence $r_1 - r_2 = 0$ implying that $r_1 = r_2$.
    Then $(q_2 - q_1)b = 0$ and $b \neq 0$ implies that $q_1 = q_2$.
\end{proof}

\begin{problem}
    \blue{Prove that every ideal in $\ZZ$ is principal.}\\
    \textbf{Hint:} If $J$ is a nonzero ideal, prove that it must contain positive integers and consider the smallest positive element of $J$.
\end{problem}

Now let's generalize to polynomials with coefficients in a field.

\begin{theorem}
    Let $K$ be a field, let $a(x), b(x) \in K[x]$ and suppose that $b(x) \neq 0$.
    Then there exists unique polynomials $q(x), r(x) \in K[x]$ such that
    \begin{enumerate}
        \item $a(x) = q(x)b(x) + r(x)$,
        \item $\deg r(x) < \deg b(x)$.
    \end{enumerate}
\end{theorem}

\begin{problem}
    \blue{Prove it!}\\
    It is crucial that the coefficients of our polynomials belong to a field $K$.
    There is exactly one crucial step where we use this hypothesis, make sure you identify where you use it.
\end{problem}

\begin{problem}
    \blue{Prove that every ideal in $K[x]$ is principal.}
\end{problem}

\end{document}
